\chapter{A Bag of Parts}
% full version: chapters/01b_neurontypes.tex

\pencilfig{1}{\pencilart{1}}{A bag of parts: almost all blue, a few red, and a handful of others.}

Imagine you are handed a bag holding eighty-six billion parts that all look alike, and asked to work out how a computer is built from them. Where would you start? An engineer would not study the shape of the parts. An engineer would ask: how many inputs does a part have, how many outputs, and what does it put out?

The neuron, as such a part is called, is surprisingly simple. It has thousands of inputs: wires from other neurons run to it. It has one output. But that output branches like a tree, and the same copy of the signal goes out to thousands of receivers. Picture a person who listens to a thousand people at once and answers with a single sentence, but in such a way that another thousand people hear it.

What is this signal? A short click. Every now and then the neuron fires a spike, a voltage blip about a thousandth of a second long. All the clicks of one neuron are the same, like the beats of one drum. So everything a neuron can say is written not in how loud it clicks but in \emph{when} it clicks: more often, less often, at once, or after a delay.

Inside, the neuron performs the simplest computation imaginable: it adds up the incoming signals and, if the sum goes over a threshold, clicks itself. Some inputs push the sum up, others push it down.

\section*{How to Sort the Bag}

When the tip of the output wire reaches a neighbor, it does not touch it. A narrow gap remains between them, and chemistry carries the signal across: the wire releases a pinch of a special substance, a transmitter, and the neighbor catches it. And here is a stroke of luck: almost every neuron releases the same substance at all the tips of its wire. So we can sort neurons by \emph{what} they release.

We will name a neuron type by the first four letters of its substance. A neuron that releases glutamate is \nt{GLUT}. One that releases gamma-aminobutyric acid is \nt{GABA}. Then come \nt{DOPA} (dopamine), \nt{SERO} (serotonin), \nt{NORA} (noradrenaline), \nt{ACET} (acetylcholine), and a few more.

Once the bag is sorted, the picture is very lopsided. Out of every thousand neurons about nine hundred sixty are \nt{GLUT}. Their signal nudges the receiver toward a click: this is a ``plus.'' About thirty-six more are \nt{GABA}; their signal pushes the receiver away from the threshold: this is a ``minus.'' In the cortex there are more of them, from a fifth to a quarter. And all the other types together come to a few neurons per hundred thousand. A drop in the ocean.

\section*{Telephone and Radio}

But this drop is built quite differently. The wire of a \nt{GLUT} or \nt{GABA} neuron leads to particular neighbors, like a telephone line: only those who are called get the signal. The tiny groups of \nt{DOPA}, \nt{SERO}, \nt{NORA}, and \nt{ACET} neurons, by contrast, work like radio stations. Their wires spread over huge regions of the brain, the substance is often released not into a gap but into the space around it, and millions of cells hear the same message at once.

The telephone lines carry data: here is a line in the picture, here is a sound of a certain pitch. The radio carries global settings: how loud everything is, whether it is worth learning, whether it is time to sleep. Any computer is the same: there are billions of memory cells and a handful of control registers. Without them the machine does not work anyway.

\section*{The Receiver Chooses the Meaning}

There is one subtlety. The substance itself does not know whether it is a ``plus'' or a ``minus.'' The molecule is a key. What happens when the key is inserted is decided by the lock, a special protein on the receiver's side. The same dopamine perks up some receivers and damps others, depending on which lock they carry. It is like one radio broadcast that different listeners take in different ways.

\section*{The ``Better Than Expected'' Signal}

The most famous radio channel is dopamine. What does it broadcast? Here is an experiment. An animal is unexpectedly given a drop of juice, and the dopamine neurons answer with a flurry of clicks. Then a lamp is lit before every juice. The animal learns that the lamp promises juice, and now the flurry comes at the lamp, and not at the juice itself. And if no juice follows the lamp, the dopamine neurons fall \emph{silent} at the moment it was expected.

\pencilfig{101}{%
% three rows: a spike raster of a DOPA neuron; lamp (amber) and juice (blue drop) above the axis
\foreach \row/\y in {1/0,2/-1.05,3/-2.1}{
  \draw[pen] (0,\y) -- (6.2,\y);
  \foreach \s in {0.3,0.75,1.1,2.15,2.6,3.05,3.45,3.9,4.45,4.95,5.4,5.85}{
    \pgfmathtruncatemacro{\gap}{(\row==3)&&(\s>3.6)&&(\s<4.7)}
    \ifnum\gap=0 \draw[pcGray, line width=0.7pt] (\s,\y) -- (\s,\y+0.3);\fi}}
% row 1: juice without a lamp -> burst at the juice
\foreach \s in {4.0,4.08,4.16,4.24,4.33}{\draw[pcAmber, line width=0.8pt] (\s,0) -- (\s,0.45);}
\draw[dotpen=pcBlue, wash=pcBlue] (4.15,0.72) circle (0.09); \draw[pcBlue, line width=0.6pt] (4.07,0.76) -- (4.15,0.92) -- (4.23,0.76);
% row 2: lamp -> burst at the lamp, nothing at the promised juice
\foreach \y in {-1.05,-2.1}{
  \foreach \s in {1.5,1.58,1.66,1.74,1.83}{\draw[pcAmber, line width=0.8pt] (\s,\y) -- (\s,\y+0.45);}
  \draw[dotpen=pcAmber, wash=pcAmber] (1.65,\y+0.72) circle (0.1);
  \foreach \a in {30,90,150}{\draw[pcAmber, line width=0.5pt] ($(1.65,\y+0.72)+(\a:0.14)$) -- ($(1.65,\y+0.72)+(\a:0.24)$);}}
\draw[dotpen=pcBlue, wash=pcBlue] (4.15,-0.33) circle (0.09); \draw[pcBlue, line width=0.6pt] (4.07,-0.29) -- (4.15,-0.13) -- (4.23,-0.29);
% row 3: lamp, no juice -> a gap where the juice was expected
\draw[pcBlue, line width=0.6pt, densely dotted] (4.15,-1.38) circle (0.09); \draw[pcBlue, line width=0.6pt, densely dotted] (4.07,-1.34) -- (4.15,-1.18) -- (4.23,-1.34);
\draw[penlite=pcRed] (3.65,-2.22) -- (3.65,-2.28) -- (4.65,-2.28) -- (4.65,-2.22);
\node[lbl, text=pcRed, anchor=north] at (4.15,-2.3) {dip};
\node[lbl, anchor=east, align=right] at (-0.1,0.15) {juice,\\no lamp};
\node[lbl, anchor=east, align=right] at (-0.1,-0.9) {lamp,\\then juice};
\node[lbl, anchor=east, align=right] at (-0.1,-1.95) {lamp,\\no juice};
}{A flurry at unexpected juice; at the lamp that promises it; and a dip at the moment the promised juice fails to come.}

One rule explains all three cases: dopamine reports not ``good'' but ``better than I expected.'' Unexpected juice is a pleasant surprise. Promised juice is no surprise at all. Promised but not given is an unpleasant one. Programs that learn from their mistakes use exactly this quantity. We will come back to it.

\section*{What to Take Away}

The brain is a huge network of parts that all look alike. Almost all of them are the ``pluses'' and ``minuses'' of the telephone network of data. And very few are radio stations that tune the whole network at once. In the next chapter we will see what can be built from ``pluses'' alone.

\fullversion{1}
